%%%%%%%%%%%%%%%%%%%%%%%%%%%%%%%%%%%%%%%%%%%%%%%%%%%%%%%%%%%%%%
%% Chapter 7: Educational Platform
%%%%%%%%%%%%%%%%%%%%%%%%%%%%%%%%%%%%%%%%%%%%%%%%%%%%%%%%%%%%%%

\section{Educational Platform}
\label{chap:educational-platform}

Alongside the annotation platform described in
Chapter~\ref{chap:annotation-platform}, which supports data collection for
model training, this project delivers an educational platform that exposes the
trained recognition and translation models
(Chapters~\ref{chap:isolated-gloss-recognition}
and~\ref{chap:sentence-level-translation}) to end users through an accessible,
user-facing interface.

\todo{PLACEHOLDER: This chapter is not yet covered by the project's existing
research and development reports and needs to be written once the educational
platform is implemented. The subsections below outline the expected structure.}

\subsection{Architecture and tech stack}
\label{sec:edu-platform-architecture}

\todo{PLACEHOLDER: Describe the educational platform's system architecture and
technology stack, following the same level of detail as
Section~\ref{sec:annotation-system-architecture} for the annotation platform
(frontend/backend frameworks, data storage, deployment infrastructure).}

\subsection{Model integration}
\label{sec:edu-platform-model-integration}

\todo{PLACEHOLDER: Describe how the trained isolated gloss recognition model
(Chapter~\ref{chap:isolated-gloss-recognition}) and/or the end-to-end pipeline
(Chapter~\ref{chap:sentence-level-translation}) are served to the platform --
model serving/inference infrastructure, latency considerations, and how
predictions are surfaced to the frontend.}

\subsection{User interface overview}
\label{sec:edu-platform-ui}

\todo{PLACEHOLDER: Describe the user-facing interface and interaction flow
(e.g.\ webcam-based live recognition, lesson/quiz structure, feedback to the
learner), including screenshots of the main views once implemented. Note: the
project repository already contains a live webcam recognition demo combining
the Transformer encoder with FAISS retrieval (commit ``live recognition
Transformer + FAISS''), which is a natural starting point for this section.}
